\documentclass[10pt,a4paper]{amsart}
\usepackage[margin=19mm]{geometry}
\usepackage{amsmath,amssymb,mathtools}
\usepackage[scr=rsfs,cal=euler]{mathalfa}
\usepackage{xfrac}
\usepackage[unicode,hidelinks,bookmarks=false]{hyperref}
\newcommand{\ie}{i.e.\ }
\newcommand{\cf}{cf.\ }
\newcommand{\resp}{respectively\ }
\newcommand{\Subsection}{\subsection}
\providecommand{\nobreakdash}{\nobreak\hbox{-}}
\setlength{\emergencystretch}{2em}
\multlinegap0pt
\usepackage{bm}
\usepackage{bbm}
\usepackage{dsfont}
\usepackage{xspace}
\usepackage{cancel}
\usepackage{mathtools}
\usepackage{amsthm}
\makeatletter
\@ifundefined{theorem}{\newtheorem{theorem}{Theorem}}{}
\@ifundefined{lemma}{\newtheorem{lemma}[theorem]{Lemma}}{}
\makeatother
\newcommand{\Hau}{{\mathcal H}} % Misura di Hausdorff
\newcommand{\Lip}{\mathop{\mathrm{Lip}}}
\newcommand{\N}{\mathbb{N}}
\newcommand{\R}{\mathbb{R}}
\newcommand{\Tr}{{\mathrm{Tr}}}
\newcommand{\cC}{\mathcal{C}}
\newcommand{\cL}{\mathcal{L}}
\newcommand{\ch}{\mathbf{1}}
\newcommand{\dd}{d}
\newcommand{\de}{\partial}
\newcommand{\difsim}{\Delta}
\newcommand{\e}{\varepsilon}
\newcommand{\loc}{\mathop{\mathrm{loc}}}
\newcommand{\pr}{P_{\Omega}}
\newcommand{\p}{P}
\newcommand{\ra}{\rightarrow}
\title[Free-Boundary Monotonicity for Almost-Minimizers of the Rela]{Free-Boundary Monotonicity for Almost-Minimizers of the Relative Perimeter}
\author{Gian Paolo Leonardi, Giacomo Vianello}
\date{Source: arXiv:2407.05039v2 (CC BY 4.0)}
\begin{document}
\maketitle

\noindent\emph{Source and licence.} Free-Boundary Monotonicity for Almost-Minimizers of the Relative Perimeter. Version arXiv:2407.05039v2 (CC BY 4.0), distributed under the Creative Commons Attribution 4.0 International licence, as recorded in the arXiv licence metadata for this exact version and shown on the abstract page https://arxiv.org/abs/2407.05039v2. Full rights notice: licenses/arxiv-2407.05039-CC-BY-4.0.txt.\par\smallskip

\noindent\textit{Editorial scope.} Counted unit: the Lemma whose source label is \texttt{lem:gaplsc} in the article (source0.tex lines 766--779, source label \texttt{lem:gaplsc}), delivered together with the article's own complete original proof of it (source0.tex lines 780--861, one \texttt{proof} environment of 82 lines). No mathematics is written, repaired or re-derived: statement and proof are the article's own text line for line. Required context retained uncounted: lem:faetrace (lines 486--491); eq:Ssymmetry (lines 576--578); eq:tildeE (lines 584--586); lem:boundS (lines 590--600); eq:A-M (lines 667--669). Every definition, display and label of those lines is retained unchanged. Omitted from this package and not redistributed: 1--485, 492--575, 579--583, 587--589, 601--666, 670--765, 862--2032. Nothing inside a retained excerpt is omitted or altered: every retained body is delivered exactly as the article's lines are written, so no omission receipt of any kind accompanies this package. Citations: the retained excerpts carry no citation command at all, so no bibliography entry of the article is transcribed and no reference is redistributed. Reference adaptation: none was needed and none was made; every reference inside the delivered unit resolves inside it, so no unresolved reference or citation is produced. Printed slips kept literal: no printed word, symbol, hypothesis or number of the counted statement or of its proof was altered. Layout and class: the article's own class is \texttt{amsart} with packages including graphicx, lmodern, babel, afterpage, color, xcolor, amsmath, amssymb, amsfonts, amsthm, enumerate, hyperref, mathrsfs, accents; the wrapper is the lane's pinned \texttt{amsart} editorial wrapper at 10pt on A4, which restates the theorem-like environments the retained text needs and adds the article's own macro definitions that the retained text needs (\texttt{\textbackslash Hau}, \texttt{\textbackslash Lip}, \texttt{\textbackslash N}, \texttt{\textbackslash R}, \texttt{\textbackslash Tr}, \texttt{\textbackslash cC}, \texttt{\textbackslash cL}, \texttt{\textbackslash ch}, \texttt{\textbackslash dd}, \texttt{\textbackslash de}, \texttt{\textbackslash difsim}, \texttt{\textbackslash e}, \texttt{\textbackslash loc}, \texttt{\textbackslash p}, \texttt{\textbackslash pr}, \texttt{\textbackslash ra}), with extra wrapper packages (bm, bbm, dsfont, xspace, cancel, mathtools). The delivered numbering is the unit's own. Assets: no retained span contains a figure environment, a graphics-inclusion command, a PDF-page inclusion, a table or a TikZ picture; no image file is copied into this package, the package declares no asset dependency and no image file is redistributed with it. Licence: version arXiv:2407.05039v2 of this article is distributed under the Creative Commons Attribution 4.0 International licence, as recorded in the arXiv licence metadata for this exact version and shown on the abstract page https://arxiv.org/abs/2407.05039v2. A full rights notice is delivered at licenses/arxiv-2407.05039-CC-BY-4.0.txt. Characters: every retained excerpt is pure ASCII, so the delivered engine, XeLaTeX with its default Latin Modern text font, reports no missing character.
\begin{lemma} \label{lem:faetrace}
Fix an open set $A \subset \R^{n}$ and a Lipschitz function $\phi:A\to \R$ of class $C^{1}$, such that $|\nabla \phi(x)|>0$ for all $x\in A$. Set $A_{r} = \phi^{-1}(-\infty,r)$. Then for all $f \in BV(A)$, for $\cL^{1}$-almost every $r\in \R$, and $\Hau^{n-1}$-almost every $x \in \de A_{r}\cap A$, we have
\begin{equation}\label{eq:fxequaltrace}
f(x) = \Tr^{+} (f, \de A_{r})(x) = \Tr^{-} (f, \de A_{r})(x) \, .
\end{equation}
\end{lemma}

\begin{equation}\label{eq:Ssymmetry}
S(x) := (x', 2 \omega(x') - x_n) \, .
\end{equation}

\begin{equation}\label{eq:tildeE}
\tilde E = E \cup (S_{\rho}(E)\setminus \Omega)\,,
\end{equation}

\begin{lemma} \label{lem:boundS}
Let $E \subset \Omega$ be a measurable set with $\ch_E \in BV_{\loc}(\Omega)$. Then, for almost all $x \in \de^\ast E$, $S$ restricted to $x+T_{x}E$ is differentiable at $x$, and we denote by $d^E S_x: T_x E \ra \R^{n}$ its differential. Moreover, if $\tilde E$ is the set defined in \eqref{eq:tildeE}, we have 
\begin{equation}\label{eq:pEtildezero}
\p(\tilde E;\de\Omega\cap \cC_{\rho,m}) = 0
\end{equation}
and, for all Borel sets $B \subset S(\cC_{\rho,m} \cap \Omega)$, 
\begin{equation} \label{eq:perim}
\p(\tilde E;B) = \p(S_{\rho}(E);B) = \int_{\de^\ast E \cap S(B)} J^E S (x) \, d \Hau^{n-1}(x) \le C\,\p(E;S(B)) \, ,
\end{equation}
where  and $C = \Lip(S)^{n-1}$.
\end{lemma}

\begin{equation}\label{eq:A-M}
\pr(E;B_{r}(x)) \leq \pr(F;B_{r}(x)) + |F \difsim E|^{\frac{n-1}{n}} \psi_{\Omega}(E;x,r) \, ,
\end{equation}

\begin{lemma}\label{lem:gaplsc}
For $j\in \N$ we let $\Omega_{j},\Omega\subset \R^{n}$ be open sets with uniformly Lipschitz boundary, such that $0\in \de\Omega$ and $\Omega_{j}\to \Omega$ locally in Hausdorff distance. Let $E_{j}, E$ be sets of locally finite perimeter, such that $E_{j}\subset \Omega_{j},\ E\subset \Omega$, and $E_{j}\to E$ in $L^{1}(B_{r_{0}})$. Finally, we assume that $E_{j}$ satisfies \eqref{eq:A-M} for all $0<r<r_{0}$ and $x=0$, and that moreover we have
\[
\lim_{r\to 0^{+}}\, \sup_{j} \psi_{\Omega_{j}}(E_{j};0,r) = 0\,.
\]
Then, $E$ satisfies \eqref{eq:A-M} for all $0<r<r_{0}$, with $\psi_{\Omega}(E;0,r) = \sup_{j} \psi_{\Omega_{j}}(E_{j};0,r)$, and 
\begin{equation}\label{LSCMinGap}
\liminf_{j} \Psi_{\Omega_{j}}(E_{j};B_{r_{0}}) \ge \Psi_{\Omega}(E;B_{r_{0}})\,.
\end{equation}
Moreover, if $\Psi_{\Omega_{j}}(E_{j};B_{r_{0}}) \to 0$ as $j\to\infty$, then for almost all $0<r<r_{0}$ we have
\begin{equation}\label{eq:perimcont}
\lim_{j \ra \infty} \p_{\Omega_{j}}(E_{j}; B_{r}) = \pr(E;B_{r}) \, , \qquad \text{for almost all $r > 0 \, .$}
\end{equation}
\end{lemma}

\begin{proof}
Let us fix $\e > 0$. Let $F \subset \Omega$ be such that $F \difsim E \subset \subset B_{r_{0}}$ and
\begin{equation}\label{eq:PsiFcontroE}
 \Psi_{\Omega}(E;B_{r_{0}}) \le \pr(E;B_{r_{0}}) - \pr(F;B_{r_{0}}) + \e\, .
\end{equation}
By Lemma \ref{lem:faetrace}, for all $j \geq 1$, for almost every $0 < r < r_{0}$ and $\Hau^{n-1}$-a.e. $x\in \de B_{r}$, we have 
\begin{equation}\label{eq:equaltraces}
\ch_{E_{j}}(x) = \Tr^{\pm}(E_{j},\de B_{r})(x)\quad \text{and} \quad  \ch_{E}(x) = \Tr^{\pm}(E,\de B_{r})(x)\,,
\end{equation}
where $\Tr^{\pm}(A,\de B_{r}) := \Tr^{\pm}(\ch_A, \de B_{r})$. By the $L^{1}_{\loc}$-convergence of $E_{j}$ to $E$, we can choose $r<r_{0}$ with the above property and the additional
\begin{equation}\label{eq:perrrzero}
\pr(E;B_{r}) \ge \pr(E;B_{r_{0}}) - \e\,,
\end{equation}
then take $j_{\e}$ large enough, such that 
\begin{align} \label{eq:difcompr}
&F \difsim E \subset \subset B_{r}\,,\\\label{eq:tracesmall}
&\int_{\de B_{r}} |\ch_{E_{j}} - \ch_{E}| \dd \Hau^{n-1} <\e\qquad \text{for $j\ge j_{\e}$} \, .
\end{align}
Let us fix $\delta>0$, define
\[
U_\delta := \big\{x=(x',x_{n})\in \cC_{r,m}:\ \omega(x')-\delta < x_{n}\le \omega(x')\big\}\,,
\]
and assume $\delta$ so small that 
\begin{equation}\label{eq:Haudeltaeps}
\Hau^{n-1}(\de B_{r}\cap U_{\delta})<\e
\end{equation}
and
\begin{equation}\label{eq:peromegadelta}
\pr(E;B_{r}) \le \p(E;B_{r}\cap(\Omega + \delta e_{n})) + \e\,.
\end{equation}
Owing to the uniform convergence of $\omega_{j}$ to $\omega$, we can select $j_\delta \geq 1$ such that $\|\omega_{j}-\omega\|_{\infty} <\delta$ for all $j\ge j_{\delta}$, then for those $j$ we define
\begin{equation*}
A_j := \Omega_{j} \cap \Omega \cap B_r \, , \qquad B_j := (\Omega_{j} \setminus \Omega) \cap B_r \,,
\end{equation*}
and observe that $B_j \subset U_\delta$. Now we set
\begin{equation} \label{eq:Fj}
F_{j} := (\tilde{F} \cap \Omega_{j} \cap B_r) \cup \big(E_{j} \cap (B_{r_{0}} \setminus B_{r})\big) \,,
\end{equation}
where $\tilde{F}=F\cup (S(F)\setminus \Omega)$ and $S$ is the symmetry through $\de \Omega$ defined in \eqref{eq:Ssymmetry} (with $\rho = r_{0}$). Thanks to \eqref{eq:difcompr}, we have $F_{j}\subset \Omega_{j}$ and $F_{j} \difsim E_{j} \subset \subset B_{r_{0}}$, which means that $F_{j}$ is a competitor for $E_{j}$ in the definition of $\Psi_{\Omega_{j}}(E_{j},B_{r_{0}})$. Moreover, by \eqref{eq:Fj}, \eqref{eq:tracesmall}, and \eqref{eq:Haudeltaeps}, we have
\begin{align} \nonumber
\p_{\Omega_{j}}(F_{j};B_{r_{0}}) & \le \p_{\Omega_{j}}(F_{j};B_{r}) + \p_{\Omega_{j}}(E_{j};B_{r_{0}} \setminus \overline{B_{r}}) \\
& \qquad + \int_{\de B_{r} \cap \Omega_{j}\cap \Omega} |\ch_{E_{j}} - \ch_{E}| d \Hau^{n-1} + \Hau^{n-1}(\de B_{r}\cap U_{\delta}) \nonumber \\
\label{eq:pF_j}
&\le \p_{\Omega_{j}}(F_{j};B_{r}) + \p_{\Omega_{j}}(E_{j};B_{r_{0}} \setminus \overline{B_{r}}) + 2\e\,.
\end{align}
Let us compute $\p(F_{j};B_{r} \cap \Omega_{j})$. By Lemma \ref{lem:boundS}, $\p(\tilde F;\de \Omega\cap \cC_{r_{0},m})=0$, hence
	    \begin{align} \label{eq:pF_jr}
	    	\p_{\Omega_{j}}(F_j;B_r) & = \p(F;A_j) + \p(S(F);B_j)\, .
	    \end{align}
        Again by Lemma \ref{lem:boundS}, up to possibly taking a smaller $\delta$, we have
\begin{equation}\label{eq:stimaFtildeBj}
\p(S(F);B_j) = \p(\tilde{F};B_{j}) \le \p(\tilde{F};U_\delta) \leq \e\,.
\end{equation}
Putting together \eqref{eq:pF_j}, \eqref{eq:pF_jr}, and \eqref{eq:stimaFtildeBj}, and taking into account \eqref{eq:equaltraces}, we obtain
\begin{align} \label{eq:EstMinGap1}
\Psi_{\Omega_{j}}(E_{j};B_{r_{0}}) & \geq \p_{\Omega_{j}}(E_{j};B_{r_{0}}) - \p_{\Omega_{j}}(F_{j};B_{r_{0}}) \\
			& \ge \p_{\Omega_{j}}(E_{j};B_{r}) - \p(F; A_{j}) -3\e\nonumber \\
			& \geq \p_{\Omega_{j}}(E_{j};B_{r}) - \pr(F;B_{r_{0}}) -3\e\, . \nonumber
		\end{align}
Now, since for all $j\ge j_{\e}$ we have $\|\omega_{j}-\omega\|_{\infty}<\delta$, we infer that
\[
B_{r}\cap (\Omega+\delta e_{n}) \subset B_{r}\cap \Omega_{j}\,.
\]
This inclusion combined with \eqref{eq:EstMinGap1} and \eqref{eq:peromegadelta} gives
\begin{align} \label{eq:EstMinGap2}
\Psi_{\Omega_{j}}(E_{j};B_{r_{0}}) & \ge \p(E_{j};B_{r} \cap (\Omega+\delta e_{n})) - \p(F;B_{r_{0}} \cap \Omega) -3\e\,,
\end{align}
so that taking the liminf as $j\to\infty$ in \eqref{eq:EstMinGap2}, and using the lower semicontinuity of the perimeter, \eqref{eq:peromegadelta}, \eqref{eq:perrrzero}, and \eqref{eq:PsiFcontroE}, we find
\begin{align*}
\liminf_{j} \Psi_{\Omega_{j}}(E_{j};B_{r_{0}}) &\ge \liminf_{j} \p(E_{j};B_{r} \cap (\Omega+\delta e_{n})) - \pr(F;B_{r_{0}}) -3\e\\
&\ge \p(E;B_{r} \cap (\Omega+\delta e_{n})) - \pr(F;B_{r_{0}}) -3\e\\
&\ge \pr(E;B_{r}) - \pr(F;B_{r_{0}}) -4\e\\
&\ge \pr(E;B_{r_{0}}) - \pr(F;B_{r_{0}}) -5\e\\
&\ge \Psi_{\Omega}(E;B_{r_{0}}) - 6\e\,.
\end{align*}
Then, the arbitrary choice of $\e$ implies \eqref{LSCMinGap}. Now, the fact that $E$ satisfies \eqref{eq:A-M} with $\psi_{\Omega}(E;0,r)$ as in the statement can be proved with the same argument used to show \eqref{LSCMinGap}, also taking into account that $|E_{j}\difsim F_{j}|\to |E\difsim F|$ as $j\to\infty$. Finally, to prove \eqref{eq:perimcont} we consider the sequence $F_{j}$ defined as before, but now with $F=E$. Choosing $\e>0$ arbitrarily, for $j$ large enough we obtain as before
\begin{align*}
\Psi_{\Omega_{j}}(E_{j};B_{r_{0}}) &\ge \p_{\Omega_{j}}(E_{j};B_{r_{0}}) - \p_{\Omega_{j}}(F_{j};B_{r_{0}})\\
&\ge \p_{\Omega_{j}}(E_{j};B_{r}) - \pr(E;B_{r}) - 3\e\,,
\end{align*}
which gives the desired conclusion.
\end{proof}

\end{document}
